\documentclass[11pt]{article}
\usepackage{fontspec}
\setmainfont{Times New Roman}
\setsansfont{Arial}
\setmonofont{Consolas}
\usepackage{amsmath,amssymb,mathtools}
\usepackage{amsthm}
\usepackage{geometry}
\usepackage{microtype}
\geometry{a4paper,margin=2.05cm}
\setlength{\parindent}{1.1em}
\setlength{\parskip}{0pt}
\newcommand{\foreign}[1]{#1}
\theoremstyle{plain}
\newtheorem{satz}{Теорема}
\newtheorem{zusatz}{Дополнение}
\emergencystretch=220em
\allowdisplaybreaks
\sloppy
\hbadness=10000
\vbadness=10000
\hfuzz=5pt
\vfuzz=12pt
\raggedbottom
\begin{document}

\newpage
\part*{Часть III\hspace{1em}Статья Капферера --- Нётер}
\addcontentsline{toc}{part}{Часть III: Статья Капферера --- Нётер}

\begin{center}
\Large\textbf{Необходимые и достаточные условия кратности}\\[0.3em]
\Large\textbf{для основной теоремы Нётер}\\[0.5em]
\large\textbf{об алгебраических функциях}\\[1em]
\normalsize\textit{Г. Капферер}\\[0.5em]
\small\textit{(с дополнением, совместно с Э. Нётер)}\\[1em]
\normalsize \foreign{Math.\ Ann.}\ \textbf{97} (1927), S.~559--567
\end{center}

\vspace{1em}

Нижеследующее представляет собой усиление основной теоремы Нётер
в форме необходимых и достаточных условий кратности\textsuperscript{1)}.
Комплекс фактов основной теоремы Нётер я, опираясь на
идеально-теоретическое понимание этой теоремы\textsuperscript{2)},
но не пользуясь здесь и далее идеально-теоретическими понятиями и
теоремами, свожу к следующим двум отдельным утверждениям\textsuperscript{3)}:

\smallskip
\textbf{Теорема Ia.} Для каждой общей нулевой точки $x = a$,
$y = b$ двух взаимно простых полиномов $\varphi$, $\psi$ от $x$, $y$
с произвольными комплексными коэффициентами существует наименьшее
целое рациональное положительное число $\varrho$ такое, что модуль
$(\varphi, \psi, \mathfrak{p}^{\varrho})$ означает в точности ту же
совокупность полиномов, что и модуль
$(\varphi, \psi, \mathfrak{p}^{\sigma})$, где
$\mathfrak{p} = (x-a, y-b)$, а $\sigma$ означает произвольное целое
рациональное число, большее $\varrho$\textsuperscript{4)}.

\smallskip
\textbf{Теорема Ib.} Если эту понятийно однозначно определенную
совокупность полиномов обозначить через $\mathfrak{q}$, соответственно
через $\mathfrak{q}_{\nu}$, когда речь идет о $\nu$-й из $s$ различных
нулевых точек, то есть о $x = a_{\nu}$, $y = b_{\nu}$, то имеет место
следующее:

Необходимо и достаточно для того, чтобы заданный полином $K$ от
$x$, $y$ обладал свойством делимости
$K \equiv 0(\varphi, \psi)$, чтобы одновременно выполнялись следующие
$s$ конгруэнций
\[
K \equiv 0(\mathfrak{q}_{1}), \quad K \equiv 0(\mathfrak{q}_{2}), \quad\ldots,\quad
K \equiv 0(\mathfrak{q}_{s}).
\]

Исторически, ср. введение к основополагающей работе Макса Нётер об
этом предмете\textsuperscript{5)}, основная теорема возникла из
попыток решить следующую фундаментальную задачу: для каждого
заданного полинома $K$ от $x$, $y$ надо уметь установить, делится ли
он на $(\varphi, \psi)$ или нет, то есть также существуют ли два
других полинома $A$ и $B$ от $x$, $y$ такие, что равенство
\[
K = A\varphi + B\psi
\]
становится тождеством по $x$, $y$.

Эта задача теоремой Нётер не решается; скорее последнюю следует
понимать лишь как, хотя и фундаментальное, сведение\textsuperscript{6)}
поставленной задачи к более простой, а именно к вопросу о
необходимых и достаточных условиях существования отдельных
конгруэнций $K \equiv 0(\mathfrak{q})$. Когда же имеет место
$K \equiv 0(\mathfrak{q})$? Только после ответа на этот последний
вопрос решена и исходная постановка задачи М. Нётер. Именно это
является целью настоящей статьи, и эта цель действительно достигается:
указываются конечное число условий кратности, практически также
легко проверяемых, которые необходимы и достаточны для
$K \equiv 0(\mathfrak{q})$. Одно достаточное условие для
$K \equiv 0(\mathfrak{q})$, которое, однако, не является одновременно
необходимым условием, известно и часто применяется и цитируется в
литературе, например в такой форме:

\smallskip
\textit{Для существования конгруэнции $K \equiv 0(\varphi, \psi)$
достаточно, чтобы каждая точка пересечения $\varphi = 0$, $\psi = 0$
входила в полином $K$ с ``достаточно'' высокой кратностью.}

\smallskip
Эта теорема содержится в теореме I как специальный случай; ведь
$\varrho$-кратная точка в $K$ равносильна
$K \equiv 0(\mathfrak{p}^{\varrho})$; тогда тем более
$K \equiv 0(\varphi, \psi, \mathfrak{p}^{\varrho})$, то есть также
$K \equiv 0(\mathfrak{q})$. Но обратное неверно.

\vspace{0.5em}
\noindent\rule{\textwidth}{0.4pt}
\vspace{0.3em}

{\footnotesize
\textsuperscript{1)} За подробным изложением я отсылаю к моей работе того же
названия в специальном выпуске отчетов Гейдельбергской академии 1927 года
``Beiträge zur Algebra'', подготовленном по инициативе A.\ Loewy,
Freiburg i.\ Br., который содержит работы разных авторов.

\textsuperscript{2)} Ср., например, B.\ L.\ van der Waerden, ``Zur Nullstellentheorie der Polynomideale''.
\foreign{Math.\ Annalen} 96 (1926), S.~203.

\textsuperscript{3)} В \S~1 цитированной работы я доказал оба утверждения
отдельно и непосредственно, то есть без идеально-теоретических понятий.

\textsuperscript{4)} Буквы $\mathfrak{p}$ и $\mathfrak{q}$ выбраны для того,
чтобы напомнить о содержательном совпадении обозначаемого с понятиями
простого идеала $\mathfrak{p}$ и первичного идеала $\mathfrak{q}$, обычными
в новейшей литературе.

\textsuperscript{5)} Max Noether, ``Über einen Satz aus der Theorie der algebraischen Funktionen'',
\foreign{Math.\ Annalen} 6 (1872).

\textsuperscript{6)} Указанное сведение приблизительно сравнимо с тем, которое
каждый знает из элементарной теории чисел: свойство делимости натурального
числа $A$ на такое же $B$ эквивалентно свойству делимости $A$ на
$p_{1}^{e_{1}}$, на $p_{2}^{e_{2}}$, \ldots, на $p_{s}^{e_{s}}$, где
$p_{1}, p_{2}, \ldots, p_{s}$ означают различные простые делители
$B = p_{1}^{e_{1}}p_{2}^{e_{2}} \cdots p_{s}^{e_{s}}$. Именно понятие
``первичный идеал'', в тексте обозначенное через $\mathfrak{q}$, надо
рассматривать как аналог степени простого числа. [Ср. E.\ Noether,
``Idealtheorie in Ringbereichen'', \foreign{Math.\ Annalen} 83 (1921), введение.]
}

\newpage
\section*{Главная теорема}

Обещанное решение задачи можно коротко охарактеризовать так:
каждому отдельному $\mathfrak{q}$ можно сопоставить систему из $t$
положительных целых рациональных чисел
$(g_{1}), (g_{2}), \ldots, (g_{t})$, а модулю $(K, \mathfrak{q})$ ---
столько же целых рациональных положительных или отрицательных чисел
$(K_{1}), (K_{2}), \ldots, (K_{t})$, так что имеет место следующая
\textbf{главная теорема}:

\begin{satz}[Главная теорема]\label{satz:kap-II}
Необходимо и достаточно для $K \equiv 0(\mathfrak{q})$, чтобы
одновременно выполнялись $t$ условий
\[
(K_{1}) \geq (g_{1}), \quad (K_{2}) \geq (g_{2}), \quad\ldots,\quad (K_{t}) \geq (g_{t}).
\]
При этом
\[
(g_{1}) + (g_{2}) + \cdots + (g_{t}) = \mu,
\]
где $\mu$ означает соответствующую $\mathfrak{q}$ кратность
пересечения, определенную результантом.
\end{satz}

В доказательстве, без ограничения общности, будем считать $x = 0$,
$y = 0$ нулевой точкой пары $(\varphi, \psi)$, соответствующей
$\mathfrak{q}$. Далее вообще, если $A(x,y)$ есть некоторый
полином\textsuperscript{7)} от $x$, $y$, символ $(A)$ будет означать
число, показывающее, сколько раз $y$ входит множителем в $A(0,y)$;
если же $A(0,y)$ тождественно обращается в нуль, то $(A)$ полагается
равным $\infty$. Числа $(g_{1}), (g_{2}), \ldots, (g_{t})$,
соответствующие $\mathfrak{q}$, определяются следующей системой
уравнений, которая образует методическую основу процедуры:
\begin{equation}\label{eq:kap-1}
\begin{aligned}
f_{1}\cdot g_{1}(0,y):y^{(g_{1})} - g_{1}\cdot f_{1}(0,y):y^{(g_{1})} &= x\cdot h_{1},\\
f_{2}\cdot g_{2}(0,y):y^{(g_{2})} - g_{2}\cdot f_{2}(0,y):y^{(g_{2})} &= x\cdot h_{2},\\
&\vdots \qquad\qquad\vdots\qquad\qquad\vdots\\
f_{r}\cdot g_{r}(0,y):y^{(g_{r})} - g_{r}\cdot f_{r}(0,y):y^{(g_{r})} &= x\cdot h_{r}.
\end{aligned}
\end{equation}
Она строится только из заданной пары полиномов $\varphi$, $\psi$:
полиномы $f_{1}, g_{1}$ определяются как тождественные с
$\varphi, \psi$, возможно после перестановки; порядок должен быть
таким, чтобы $(f_{1}) \geq (g_{1})$. При этой фиксации, и только
тогда, $h_{1}$ будет целорациональным.

Полиномы $f_{2}, g_{2}$ определяются как тождественные с
$h_{1}, g_{1}$ при дополнительном условии
$(f_{2}) \geq (g_{2})$. В общем случае пара полиномов
$f_{r}, g_{r}$ определяется как тождественная с $h_{r-1}, g_{r-1}$
при дополнительном условии $(f_{r}) \geq (g_{r})$.

Систему уравнений \eqref{eq:kap-1} действительно можно продолжать
бесконечно указанным способом; но главная теорема требует построить
только первые $t$ уравнений системы, где $t$ определяется следующим
фактом:

\smallskip
\textbf{Лемма I.} Существует натуральное число $t$ [причем
$t \leq \mu$, где под $\mu$ понимается кратность пересечения, с
которой точка $x = 0$, $y = 0$ входит в $\varphi = 0$, $\psi = 0$]
такое, что первые $t$ чисел, определенных посредством
\eqref{eq:kap-1}, а именно $(g_{1}), (g_{2}), \ldots, (g_{t})$, все
больше нуля, тогда как одновременно $(g_{t+1})$ и все следующие
равны нулю.

Действительно, с учетом теоремы о произведении результантов кратность
пересечения в точке $x = 0$, $y = 0$ для $f_{r} = 0$, $g_{r} = 0$
равна
\[
\mu - (g_{1}) - (g_{2}) - \cdots - (g_{r}),
\]
откуда следует обрыв не более чем через $\mu$ шагов. Отсюда
одновременно получается факт
\begin{equation}\label{eq:kap-2}
\mu = (g_{1}) + (g_{2}) + \cdots + (g_{t}),
\end{equation}
формула, замечательная тем, что в общем случае она дает рациональный
и практически легко выполнимый метод определения кратности
пересечения, как только известна соответствующая нулевая точка.

Когда числа $(g_{i})$ уже известны, числа $(K_{i})$ определяются
следующей системой уравнений, соответствующей \eqref{eq:kap-1}:
\begin{equation}\label{eq:kap-3}
\begin{aligned}
K_{1}\cdot g_{1}(0,y):y^{(g_{1})} - g_{1}\cdot K_{1}(0,y):y^{(g_{1})} &= x\cdot K_{2},\\
K_{2}\cdot g_{2}(0,y):y^{(g_{2})} - g_{2}\cdot K_{2}(0,y):y^{(g_{2})} &= x\cdot K_{3},\\
&\vdots \qquad\qquad\vdots\qquad\qquad\vdots\\
K_{r}\cdot g_{r}(0,y):y^{(g_{r})} - g_{r}\cdot K_{r}(0,y):y^{(g_{r})} &= x\cdot K_{r+1}.
\end{aligned}
\end{equation}
$K_{1}$ определяется как тождественный с $K$. Очевидно, $K_{2}$ будет
целорациональным тогда и только тогда, когда $(K_{1}) \geq (g_{1})$.
Если последнее выполнено, далее следует: $K_{3}$ будет
целорациональным тогда и только тогда, когда $(K_{2}) \geq (g_{2})$.
В общем случае, если $K_{r}$ уже целорационален, $K_{r+1}$ будет
целорациональным тогда и только тогда, когда $(K_{r}) \geq (g_{r})$.

Доказательство главной теоремы (теоремы~\ref{satz:kap-II}) теперь
основано на следующей

\smallskip
\textbf{Лемме II.} Если положить
$(f_{i}, g_{i}, \mathfrak{p}^{e})$ равным $\mathfrak{q}^{(i)}$, то для
каждого $i$ из ряда целых чисел $1, 2, \ldots$ \textit{in inf.}
имеет место утверждение:

Для конгруэнции
\[
K_{i} \equiv 0(\mathfrak{q}^{(i)})
\]
необходимо и достаточно одновременное выполнение двух условий:
\[
(K_{i}) \geq (g_{i}) \qquad \text{и} \qquad K_{i+1} \equiv 0(\mathfrak{q}^{(i+1)}).
\]

Доказательство\textsuperscript{8)} леммы получается комбинацией
систем уравнений \eqref{eq:kap-1} и \eqref{eq:kap-3}. Если учесть,
что $g_{t+1}$ согласно лемме I уже не обращается в нуль в нулевой
точке, то есть что модуль $\mathfrak{q}^{(t+1)}$ состоит из всех
полиномов (является единичным идеалом), и потому
$K_{t+1} \equiv 0(\mathfrak{q}^{(t+1)})$ становится тривиальным, то
$t$-кратное применение леммы II дает доказательство главной
теоремы.

\smallskip
Без доказательств, доказательства находятся в моей цитированной
выше работе, я добавляю еще два дополнения к главной теореме:

\begin{zusatz}[I]
Число $t$ главной теоремы тождественно с наименьшим\textsuperscript{9)}
положительным числом $\nu$ такого рода, что
$x^{\nu} \equiv 0(\mathfrak{q})$.
\end{zusatz}

\begin{zusatz}[II]
Системы уравнений \eqref{eq:kap-1} и \eqref{eq:kap-3}, а также
возникающие из них перестановкой $x$ и $y$, при подходящей
комбинации дают рациональный метод точного численного определения
``принадлежащего'' $\mathfrak{q}$ показателя $\varrho$, то есть того
минимального показателя $\varrho$ из теоремы Ia.
\end{zusatz}

Э.\ Бертини\textsuperscript{10)} дал общую формулу для $\varrho$.
Что эта формула, однако, не всегда дает минимальное значение
$\varrho$, показано в цитированной работе на примере
$\varphi = x^{3} + y^{4}$, $\psi = x^{2} + y^{3}$. Наконец заметим,
что и сама формула Бертини может быть по-новому обоснована нашей
главной теоремой.

\vspace{0.5em}
\noindent\rule{\textwidth}{0.4pt}
\vspace{0.3em}

{\footnotesize
\textsuperscript{7)} Если для $A(x,y)$ допустить также дробно-рациональные
функции $C(x,y):D(x,y)$, то $(A) = (C) - (D)$; в последнем случае
$(A)$ поэтому может принимать и отрицательные значения.

\textsuperscript{8)} Ср. также доказательство в замечании I ``Дополнения''.

\textsuperscript{9)} Здесь, следовательно, получаются точные показатели в
отличие от оценок фрейлейн Грете Герман [``Die Frage der endlich vielen Schritte
in der Theorie der Polynomideale'' (с использованием посмертных теорем
К. Гентцельта), \foreign{Math.\ Annalen} 95 (1926)], где, правда, при
гораздо более общей постановке задачи, указываются только верхние границы.

\textsuperscript{10)} E.\ Bertini, \foreign{Math.\ Annalen} 34 (1889).
}

\newpage
\section*{Теорема III и дополнение}

Интересная специализация главной теоремы следует из ограничительного
предположения, что нулевая точка соответствующего
$\mathfrak{q} = (\varphi, \psi, \mathfrak{p}^{e})$ является
простой точкой по крайней мере на одной из двух кривых
$\varphi = 0$, $\psi = 0$. Если, например, она является простой
точкой на $\psi$, то можно показать, что число $t$ главной теоремы
точно равно $\mu$ и что $t$ неравенств главной теоремы можно
заменить одним-единственным условием кратности. Это приводит к
следующей общей теореме:

\begin{satz}[Теорема III]\label{satz:kap-III}
Если все точки пересечения $\varphi = 0$, $\psi = 0$ являются
простыми точками на $\psi$, то, при условии что $K$ взаимно прост с
$\psi$, для существования конгруэнции $K \equiv 0(\varphi, \psi)$
необходимо и достаточно, чтобы $K$ обращался в нуль во всех этих
точках, причем так, чтобы кратность пересечения в паре кривых
$K = 0$, $\psi = 0$ в каждой из них была не меньше, чем в паре
кривых $\varphi = 0$, $\psi = 0$.
\end{satz}

Эта теорема III замечательна еще по особой причине: требуемое ею
условие кратности является таким, что его выполнение или невыполнение
можно установить конечным числом дифференцирований, независимо от
формулы \eqref{eq:kap-2}, основанной на \eqref{eq:kap-1}. Это сразу
становится ясным, если процитировать следующую ``теорему'', которую
я сформулировал в 1923 году\textsuperscript{11)}:

\smallskip
\textit{Если $A$ и $B$ --- какие-либо два взаимно простых полинома
от $x$, $y$, а $P$ --- их общая точка, являющаяся простой точкой для
$B$, то кратность этой точки как точки пересечения тогда и только
тогда равна числу $\mu$, когда кратность той же точки в паре кривых
$A' = 0$, $B = 0$ в точности равна $\mu - 1$; при этом}
\[
A' = \frac{\partial A}{\partial x} \cdot \frac{\partial B}{\partial y} -
\frac{\partial A}{\partial y} \cdot \frac{\partial B}{\partial x}.
\]

Наконец заметим еще, что главная теорема
(теорема~\ref{satz:kap-II}) с небольшой модификацией может быть
перенесена и на однородные полиномы от 3 переменных $x$, $y$, $z$.
Такой перенос в указанном выше специальном случае удается особенно
элегантно. А именно можно показать:

\begin{zusatz}[к теореме III]
Теорема~\ref{satz:kap-III} остается совершенно неизменной, даже по
формулировке, если под $\varphi$, $\psi$, $K$ понимать однородные
полиномы от трех переменных\textsuperscript{12)}.
\end{zusatz}

\vspace{0.5em}
\noindent\rule{\textwidth}{0.4pt}
\vspace{0.3em}

{\footnotesize
\textsuperscript{11)} H.\ Kapferer, ``Über die Multiplizität der Schnittpunkte von zwei algebraischen
Kurven''. \foreign{Jahresber.\ der D.\ M.-V.} 32 (1923).

\textsuperscript{12)} Вопрос об условиях существования однородной конгруэнции
от трех переменных
$K(x,y,z) \equiv 0(\varphi(x,y,z), \psi(x,y,z))$, хотя и при сильно
специализированных предположениях, является также центральным
вопросом исследования A.\ Ostrowski, ``Über die Maxwellsche Erzeugung der
Kugelfunktionen'', \foreign{Jahresb.\ der D.M.-V.} 33 (1925). В этой
заметке предполагается: $\varphi$ распадается только на линейные
множители; $\varphi$ и $K$ имеют один и тот же порядок;
$\psi = x^{2} + y^{2} + z^{2}$. Поскольку здесь $\psi$ совершенно
свободен от особенностей, главную теорему заметки Островского можно
рассматривать как специальный случай нашей теоремы III.
}

\newpage
\section*{Дополнение, совместно с Э. Нётер\textsuperscript{13)}}

Система уравнений \eqref{eq:kap-1} одновременно дает рациональный
метод определения полной системы вычетов по $\mathfrak{q}$, под чем
следует понимать полную систему представителей конечного числа
линейно независимых относительно области коэффициентов $P$ классов
вычетов modulo $\mathfrak{q}$. Тем самым отдельные шаги в
доказательстве главной теоремы становятся очень прозрачными; вместе
с тем в качестве побочного результата получается факт, что
результантная кратность $\mu$ совпадает с длиной идеала
$\mathfrak{q}$;

\vspace{0.5em}
\noindent\rule{\textwidth}{0.4pt}
\vspace{0.3em}

{\footnotesize
\textsuperscript{13)} Идеально-теоретическое истолкование принадлежит
Э.\ Нётер, а лежащее в основе доказательство (5) и (6) --- Г.\ Капфереру.

\textsuperscript{14)} Поле $P$ без ограничения общности предполагается
алгебраически замкнутым; в этом объеме действительно справедливы все
предыдущие теоремы, за исключением условия дифференцирования. Длина
$\mathfrak{q}$ определяется также как длина композиционного ряда из
идеалов, идущего от $\mathfrak{p}$ к $\mathfrak{q}$. Ср. E.\ Noether,
\foreign{Jahresber.\ d.\ d.\ Math.-Ver.} 34 (1925), S.~101
(косая пагинация); подробнее у H.\ Grell, \foreign{Math.\ Annalen} 97
(1927), S.~529.

Имеющиеся в литературе доказательства совпадения обоих чисел
кратности (для $n$ переменных) очень сложны; самым простым является
еще не опубликованное посмертное доказательство Гентцельта.
}

\vspace{1em}
то есть с рангом кольца классов вычетов или с числом линейно
независимых относительно $P$ классов вычетов\textsuperscript{14)}.

Полная система вычетов, в обозначениях \eqref{eq:kap-1} и леммы II,
характеризуется следующей

\begin{satz}[Теорема IV]\label{satz:kap-IV}
Полная система вычетов по $\mathfrak{q}$ задается полиномами
\[
h_{t},\ h_{t}y,\ldots,h_{t}y^{(g_{t})-1};\ \ldots;\
h_{i},\ h_{i}y,\ldots,h_{i}y^{(g_{i})-1};\ \ldots;\
h_{1},\ h_{1}y,\ldots,h_{1}y^{(g_{1})-1}.
\]
При этом каждый раз $h_{i}, h_{i}y, \ldots, h_{i}y^{(g_{i})-1}$ дает
полную систему вычетов идеала $\mathfrak{q}^{(i+1)}$ по модулю
$\mathfrak{q}^{(i)}$.
\end{satz}

Из определения непосредственно следует:
$\mathfrak{q}^{(i+1)} = (\mathfrak{q}^{(i)}, h_{i})$, значит
$\mathfrak{q}^{(i)} \equiv 0(\mathfrak{q}^{(i+1)})$.
Теорема~\ref{satz:kap-IV} будет доказана, если показать:
\begin{equation}\label{eq:kap-4a}
h_{i}x \equiv 0(\mathfrak{q}^{(i)}); \qquad
h_{i}y^{(g_{i})} \equiv 0(\mathfrak{q}^{(i)})
\end{equation}
и
\begin{equation}\label{eq:kap-5}
\text{из } h_{i}F(y) \equiv 0(\mathfrak{q}^{(i)})
\text{ следует } F(y) \equiv 0(y^{(g_{i})}).
\end{equation}

Соотношения \eqref{eq:kap-4a} утверждают, что линейные комбинации
$h_{i}, h_{i}y, \ldots, h_{i}y^{(g_{i})-1}$ исчерпывают классы
вычетов идеала $\mathfrak{q}^{(i+1)}$ по модулю
$\mathfrak{q}^{(i)}$, тогда как \eqref{eq:kap-5} выражает линейную
независимость. Так как $\mathfrak{q}^{(t+1)}$ становится единичным
идеалом, что теперь следует из конечной длины $\mathfrak{q}$, отсюда
получается доказательство теоремы~\ref{satz:kap-IV} и одновременно
факт, что длина $\mathfrak{q}$ равна
$(g_{1}) + \cdots + (g_{t})$, то есть по \eqref{eq:kap-2} совпадает
с результантной кратностью.

Первое из соотношений \eqref{eq:kap-4a} считывается из
\eqref{eq:kap-1}; второе получается следующим образом:

Имеем $h_{i}g_{i} \equiv 0(\mathfrak{q}^{(i)})$; следовательно, ввиду
первого соотношения \eqref{eq:kap-4a}, также
$h_{i}g_{i}(0,y) \equiv 0(\mathfrak{q}^{(i)})$, или
$h_{i}y^{(g_{i})}(g_{i}(0,y):y^{(g_{i})}) \equiv 0(\mathfrak{q}^{(i)})$,
откуда $h_{i}y^{(g_{i})} \equiv 0(\mathfrak{q}^{(i)})$
вследствие $g_{i}(0,y):y^{(g_{i})} \not\equiv 0(y)$.

Перед доказательством \eqref{eq:kap-5} заметим то, что
непосредственно считывается из \eqref{eq:kap-1}: $f_{i}$ и $g_{i}$
могут иметь общим делителем только не делящийся на $y$ полином
$d_{i}(y)$ от $y$: $f_{i} = d_{i}(y)\bar{f}_{i}$,
$g_{i} = d_{i}(y)\bar{g}_{i}$, причем $\bar{f}_{i}$ и
$\bar{g}_{i}$ взаимно просты. Так как
$d_{i}(y) \not\equiv 0(\mathfrak{p})$, то
$\mathfrak{q}^{(i)}$ также равно
$(\bar{f}_{i}, \bar{g}_{i}, \mathfrak{p}^{e})$; при этом оно равно
$(\bar{f}_{i}, \bar{g}_{i}, \mathfrak{p}^{\sigma})$ для каждого
$\sigma \geq \varrho$, будучи делителем $\mathfrak{q}$ при каждом
таком $\sigma$; следовательно, $\mathfrak{q}^{(i)}$ есть
соответствующая нулевой точке первичная компонента
$(\bar{f}_{i}, \bar{g}_{i})$.

Пусть $\gamma_{1}, \delta_{1}; \ldots; \gamma_{r}, \delta_{r}$ ---
отличные от $x = 0$, $y = 0$ нулевые точки
$(\bar{f}_{i}, \bar{g}_{i})$, и пусть
$\tau_{1}, \ldots, \tau_{r}$ --- показатели соответствующих
первичных компонент $(\bar{f}_{i}, \bar{g}_{i})$; далее в произведении
\[
P = ((x-\gamma_{1}) + u(y-\delta_{1}))^{\tau_{1}} \cdots
((x-\gamma_{r}) + u(y-\delta_{r}))^{\tau_{r}}
\]
неопределенную $u$ специализируем в элемент из $P$ так, чтобы
сохранилось $P \not\equiv 0(\mathfrak{p})$. Положив $G(x,y)$ равным
произведению специализированного $P$ с $d_{i}(y)$, получаем
$G(x,y) \not\equiv 0(\mathfrak{p})$ и
$\mathfrak{q}^{(i)} \cdot G(x,y) \equiv 0(f_{i}, g_{i})$.

Из предположения $h_{i}F(y) \equiv 0(\mathfrak{q}^{(i)})$ следует
$h_{i}F(y)G(x,y) \equiv 0(f_{i}, g_{i})$, или, вместе с
\eqref{eq:kap-1}, в форме тождеств:
\[
h_{i}F(y)G = Af_{i} + Bg_{i}; \qquad
h_{i}x = (g_{i}(0,y):y^{(g_{i})})\cdot f_{i}
       - (f_{i}(0,y):y^{(g_{i})})g_{i}.
\]
А это дает
\[
\bar{f}_{i}(xA - F(y)H) \equiv 0(\bar{g}_{i}); \qquad
H = G(x,y) \cdot g_{i}(0,y):y^{(g_{i})} \not\equiv 0(\mathfrak{p}).
\]
Из взаимной простоты $\bar{f}_{i}$ и $\bar{g}_{i}$ далее следует
\[
xA - F(y)H \equiv 0(\bar{g}_{i});
\]
следовательно,
$F(y)H \equiv 0(\bar{g}_{i}, x)
\equiv 0(\bar{g}_{i}, x, \mathfrak{p}^{\sigma})
= (y^{(g_{i})}, x)$, и потому $F(y) \equiv 0(y^{(g_{i})})$.
Тем самым доказано \eqref{eq:kap-5}, а значит и
теорема~\ref{satz:kap-IV}.

\smallskip
\textbf{Замечание I.} Соотношения \eqref{eq:kap-4a} и
\eqref{eq:kap-5} вместе с леммой II можно записать как взаимно
обратные соотношения между частными идеалов\textsuperscript{15)};
при этом доказательство второго соотношения одновременно дает
доказательство леммы II:
\begin{equation}\label{eq:kap-6}
\begin{aligned}
\mathfrak{q}^{(i)} : \mathfrak{q}^{(i+1)}
  &= (\mathfrak{q}^{(i)}, x) = (y^{(g_{i})}, x);\\
\mathfrak{q}^{(i)} : (\mathfrak{q}^{(i)}, x)
  &= \mathfrak{q}^{(i+1)}.\qquad\text{\textsuperscript{16)}}
\end{aligned}
\end{equation}

Первое соотношение является объединением \eqref{eq:kap-4a} и
\eqref{eq:kap-5}; по поводу второго надо заметить:
$x \cdot \mathfrak{q}^{(i+1)} \equiv 0(\mathfrak{q}^{(i)})$ было
доказано в \eqref{eq:kap-4a}; обратное получается аналогично
\eqref{eq:kap-5}. Из $xM \equiv 0(\mathfrak{q}^{(i)})$ как выше
следует:
\[
xMG \equiv 0(f_{i}, g_{i}); \quad \text{следовательно}\quad
xMG \cdot g_{i}(0,y):y^{(g_{i})} \equiv 0(xh_{i}, g_{i}).
\]
Так как $g_{i} \not\equiv 0(x)$ по определению, именно
$(f_{i}) \geq (g_{i})$, отсюда получается
\[
M \cdot (G \cdot g_{i}(0,y):y^{(g_{i})})
\equiv 0(h_{i}, g_{i}) \equiv 0(\mathfrak{q}^{(i+1)}),
\quad \text{то есть}\quad M \equiv 0(\mathfrak{q}^{(i+1)}).
\]

Необходимость условия леммы II теперь следует так:
$K_{i} \equiv 0(\mathfrak{q}^{(i)})$ дает
$K_{i}(0,y) \equiv 0(y^{(g_{i})}, x)$; значит
$(K_{i}) \geq (g_{i})$, и тем самым, по \eqref{eq:kap-3},
\[
x \cdot K_{i+1} \equiv 0(\mathfrak{q}^{(i)}); \quad
\text{следовательно, по второму соотношению \eqref{eq:kap-6},}\quad
K_{i+1} \equiv 0(\mathfrak{q}^{(i+1)}).
\]

Достаточность условия считывается непосредственно; конечное
многократное применение леммы II дает главную теорему. Точно так же
итерация второго соотношения \eqref{eq:kap-6} дает дополнение I, а
именно $\mathfrak{q} : x^{t} = \mathfrak{q}^{(t+1)}$, то есть
единичный идеал; причем $t$ является наименьшей такой степенью, так
как $\mathfrak{q} : x^{t-1} = \mathfrak{q}^{(t)}$.

\textbf{Замечание II.} Для полинома $K$, не делящегося на
$\mathfrak{q}$, можно еще утверждать: если
$K \equiv 0(\mathfrak{q}^{(i+1)})$,
$\not\equiv 0(\mathfrak{q}^{(i)})$, то
$(K, \mathfrak{q}^{(i)}) = (h_{i}y^{(K)}, \mathfrak{q}^{(i)})$.
Ведь по теореме~\ref{satz:kap-IV}
$K \equiv h_{i}y^{\lambda}c(y)$ modulo
$(\mathfrak{q}^{(i)})$ с $c(y) \not\equiv 0(y)$; следовательно,
\[
(K, \mathfrak{q}^{(i)}) = (h_{i}y^{\lambda}, \mathfrak{q}^{(i)}).
\]
При этом $y^{(g_{i})-\lambda}K \equiv 0(\mathfrak{q}^{(i)})$, и именно
это по \eqref{eq:kap-5} является наименьшей такой степенью. По лемме II
эта наименьшая степень равна $(g_{i}) - (K)$; значит
$\lambda$ равно $(K)$.

\vfill
\begin{flushright}
(Поступила 24.\ 10.\ 1926.)
\end{flushright}

\vspace{0.5em}
\noindent\rule{\textwidth}{0.4pt}
\vspace{0.3em}

{\footnotesize
\textsuperscript{15)} Под частным $\mathfrak{a}:\mathfrak{b}$, как известно,
понимают совокупность полиномов $c(x,y)$ таких, что
$\mathfrak{b}c \equiv 0(\mathfrak{a})$.

\textsuperscript{16)} Из этих двух взаимно обратных соотношений одно
влечет другое по общей теории неприводимых идеалов; ср. Macaulay,
\textit{Modular Systems}, Nr.\,72 и 73. \foreign{Cambridge Tracts} 19 (1916).
}

\end{document}
